\section{Computational Evidence and the Cost of Certification}\label{sec:study59}
The experiments address three separate questions: how an exact exception-count algorithm scales; whether an inexpensive price bound already supports an implementable book; and how much end-to-end certification costs beyond optimization. All inputs are constructed and disclosed. The study does not establish adoption, calibration to a provider, or a universally superior routing rule. Tables are collected after the references; complete input, source, execution, and proof records accompany the paper.

\subsection{Preserved evidence and an expanded frozen design}
The earlier study comprises 90 distinct model specifications, 270 budget-tagged cases, and 1,260 method requests. In its three-second standard-tariff panel, the tariff, price, deficit, and enumeration methods attain an independently checked rational target in 30, 29, 12, and 19 of 30 requests, respectively. SCIP attains its numerical target in 30 of 30. The corresponding exceptional-tariff counts are the same. In the curved panel, the counts are 28 for price, 10 for deficit, and 18 for enumeration; the tariff method is mathematically inapplicable, and SCIP reports 27 numerical successes. These results, the preceding resource-grid experiments, and all unfavorable outcomes are preserved unchanged. A favorable guarantee is not interpreted as evidence that the earlier deficit implementation was fastest.

The expanded design contains 44 distinct physical/tariff specifications, 88 budget-tagged cases, and 405 declared requests. Its controlled scaling panel crosses $(N,k,m)=(17,12,8),(65,32,16),(129,48,32)$ with $d\in\{0,1,2,4,6,8,10,12\}$ and one- and six-second total allowances. Each cell runs tariff, price, enumeration, SCIP, and a specified hybrid. Eight further tariff requests use 30 seconds on the 17-command cases. Three 13-exception requests test the engineering guard. A separate phase-split panel varies the optimization allocation between 35\% and 80\% of a three-second allowance. Five numerical-encoding cases use probability constructions based on 16, 64, 256, 1,024, and 4,096 bits, for tariff and price under 15 seconds. These are controlled cells, not random replicates.

A separately specified adversarial panel has twelve cases at $(N,k,m)=(19,10,6)$ and $(41,20,12)$. It varies exception density, a low accepted promise, many small fee deviations, positive service costs and curvature, eligibility geometry, and alternating fee levels. It does not reuse the earlier seeds or the hardness reduction. Its purpose is broader stress coverage, not a field-data claim. It adds the robust-tariff and deficit methods, giving 84 six-second requests. The full formulas and rational inputs are frozen before the expanded execution.

\subsection{The exception parameter in time, memory, and proof size}
Table~\ref{tab:scaling63} reports the controlled success counts with their full denominators. At six seconds the exact tariff method certifies all eight 17-command cases, five of eight 65-command cases, and three of eight 129-command cases. The largest certified exception counts at those sizes are 12, 6, and 2, respectively. Thus the parameterized result is reflected in an increasingly costly exception loop, not an assertion that every value up to the implementation guard is inexpensive. All eight long-allowance tariff requests certify the target.

Table~\ref{tab:d63} separates optimization, serialization, checking, resident memory, and uncompressed proof size for every exception count at the fixed 17-command scale. Increasing $d$ from zero to twelve changes the optimization time from 0.005 to 0.819 seconds and the independent check from 0.209 to 1.116 seconds; proof size increases from 4,143 to 3,964,865 bytes. At 65 commands and eight exceptions, an optimizer can finish while the parent deadline prevents completion of verification. That request remains a failure of the six-second end-to-end target. Proof files in the entire expanded panel reach 5,572,345 uncompressed bytes. Reported phase costs condition on the phase actually being reached, and optimizer and checker peak memory are distinct measurements.

The tariff algorithm certifies all five numerical-encoding cases, as does price, without a resource grid. These finite cases test exact rational arithmetic with large encodings; they do not remove bit-operation costs from the theorem. Enumeration attains none of the 48 controlled scaling targets within these allowances. The result illustrates the difference between the exception-count compression and book enumeration at these sizes; it is not an empirical proof of an asymptotic lower bound. The adversarial panel has rational successes of 4/12 for tariff, 7/12 for robust tariff, 12/12 for price, 8/12 for hybrid, and 0/12 for both enumeration and deficit. SCIP reports 12/12 numerical successes, with only original-feasible lower policies independently certified rationally.

\subsection{Same-book support: an independent exhaustive diagnostic}\label{sec:root63}
The expanded large-instance price logs contain 69 requests with an observed root probe, and none of the particular returned active books supports the promise at its recorded probe price. This does \emph{not} prove that the whole active-book family lacks support: ties and a nonminimizing probe price matter. We therefore add a separate, fully enumerated diagnostic rather than treat absence in a trace as a structural conclusion.

This diagnostic uses four equally likely caps $(0,1/3,2/3,1)$, seven catalog commands $i/6$, two ceiling patterns, allowances two and four, three promise fractions of the mean cap, and four fee patterns. The complete factorial design has 48 cases, frozen before its own execution. For every feasible book, an independent routine constructs the three vertices of its exact linear/free-service value function: its first command, terminal capacity, and mean cap. The upper concave hull of the union gives the exact minimum price bound at the promise. All active books and their maximizing intervals are then enumerated, independently of the price-path recurrence. Finally, the price implementation is run and its certificate is checked against the external input.

Across 1,176 enumerated feasible books, 20 cases have same-book support and 28 have a strictly positive minimum root gap; the largest gap is $1/20$. The equality characterization and the distance bound of Theorem~\ref{thm:pricebridge60} hold in every checked case. Table~\ref{tab:root63} contrasts observed branching and checking work by this independently established classification. All 48 final price certificates are exact. The comparison identifies a mechanism on this balanced small panel; it is not a causal estimate or a population frequency, and its local timings are not pooled with the earlier execution environment.

\subsection{A tested hybrid, and the limits of a routing claim}
The expanded hybrid spends at most one quarter of its optimization allowance on a one-node, three-price probe. If the certified gap is small enough it returns that proof. Otherwise, in the linear/free-service regime it tries the near-standard projection with an exception guard of four when the projected distortion meets the tolerance; failing that test, it restarts price search. The restart cost is charged rather than described as a free warm start. This rule was executed, not merely proposed.

Its 27/48 controlled scaling successes compare with 25/48 for tariff and 24/48 for price. At the largest size and six seconds it certifies three of eight cases while price certifies none; at one second it supplies no advantage at either larger size. In the adversarial panel price certifies all twelve cases and the hybrid eight. These outcomes do not validate a generally optimal method-selection hierarchy. They support reporting the observable inputs---physical applicability, exact or approximate exception count, and an available support witness---together with the full cost of checking the resulting decision.

The structural regression suite additionally checks 198 robust instances against 198 exhaustive fee projections and 6,623 book-error comparisons, rejects 108 corrupt certificates, and exercises exact exception counts through twelve and the guard at thirteen. The electronic companion explains the external-input serialization repair and the complete independent replay. Neither successful post-study replay nor a new packaging build retroactively converts an interrupted request into an on-time success.

\section{Conclusion}\label{sec:conclusion63}
The finite-catalog model admits an exact resource-path representation, but representation alone does not determine tractability. With common linear rewards and free service, arbitrary command-specific charges retain parameterized hardness in the command allowance and cap count, while a small number of exceptions to a standard charge permits an exact fixed-parameter algorithm. The uniform-fee frontier gives explicit installation decisions, and the perturbation certificate extends that algorithm to near-standard tariffs with a bound on the value lost by an original-feasible policy. The remaining conservation, approximation, heterogeneous-reward, prototype, and price results specify complementary regimes rather than a universal algorithm.

Computational certification has its own costs. Larger catalogs, larger exception counts, and independent checking can exhaust a deadline even when an optimizer produces a candidate. The expanded controlled and adversarial records, and the independent root-support diagnostic, make these distinctions observable. A practical implementation should preserve both the policy and its precise scope of guarantee, including the model input, fee approximation, resource equality, and time allowed for verification.

\paragraph{Code and data.} The accompanying archive contains the complete constructed inputs, unchanged historical execution sources and records, current independent checkers, reconstruction instructions, and all retained derivations. No proprietary operational data are used.
